\section{Introduction}

The expansion of the Internet of Things (IoT), together with the deployment of data-intensive applications over 5G and emerging 6G networks, has increased the demand for low-latency, high-throughput, and reliable data access at the network edge. Applications such as smart city monitoring, autonomous transportation, industrial automation, augmented reality, and large-scale sensor networks continuously generate and consume delay-sensitive, making centralized cloud processing inefficient due to latency, backhaul congestion, and scalability limitations.

Edge computing address this challenges by bringing computation and storage resources closer to end users. Among various edge computing mechanisms, content caching at the edge plays a critical role in reducing end-to-end latency and alleviating network load by storing frequently accessed data near consumers. Prior work has demonstrated that cooperative and distributed caching can significantly improve network performance in wireless systems~\cite{Shanmugam2013} and that proactive caching can further enhance quality of service under heavy and bursty traffic conditions~\cite{Bastug2014}.

However, IoT edge environments differ fundamentally from traditional web or content delivery network (CDN) caching~\cite{Podlipnig2003}. IoT systems are characterized by high user mobility, heterogeneous devices, non-stationary demand, and strict resource constraints. In such settings, data access patterns are often influenced by spatial locality, temporal dynamics, and user movement across edge nodes. As a result, caching decisions made independently at each node may lead to redundant content replication, inefficient storage utilization, and degraded system-wide performance.

To mitigate these limitations, prior research has explored cooperative caching mechanisms that allow neighboring edge nodes to coordinate content placement~\cite{Shanmugam2013}. More broadly, recent system level work has emphasized that future intelligent edge systems require stronger collaboration across distributed vehicles, roadside units, base stations, and cloud infrastructure in order to improve latency, scalability, and reliability. Bai et al. \cite{bai2025collaborative}described this direction as collaborative edge intelligence, where distributed edge resources are pooled to support low-latency and large-scale intelligent services.However, if this improves over local caching, existing approaches assume static network topology rely on predefined neighborhood structures. Even recent machine learning caching strategies~\cite{Shuja2021JNCA} focus on local demand prediction or reinforcement learning at individual nodes, treating topology implicitly or as a secondary feature.

In parallel, graph-based learning methods particularly Graph Neural Networks (GNNs) have demonstrated strong capability in modeling relational dependencies in communication networks, including routing and resource allocation~\cite{Rusek2019,Jiang2022Survey}. These methods naturally encode structural information and generalize across varying network topologies. However, their use in cooperative caching for IoT edge environments, especially under mobility driven dynamic topology, remains limited.

This work is motivated by the observation that user mobility implicitly defines a time varying interaction graph among edge nodes. Nodes that frequently exchange users tend to have correlated demand patterns and leveraging this structure through learning-based approaches may allow for more effective and coordinated caching decisions. Therefore, understanding how topological context influences cooperative caching behavior is essential for designing scalable and adaptive edge systems in future IoT and 6G networks.